% !TEX program = xelatex
% Компилација: xelatex -> xelatex (двапати за референци/содржина)
\documentclass[11pt,a4paper]{article}

%---------------------------------------------------------------
% Кодирање, јазик и фонтови (XeLaTeX + polyglossia + Unicode фонт)
%---------------------------------------------------------------
\usepackage{fontspec}
\usepackage{polyglossia}
\setdefaultlanguage{macedonian}
\setotherlanguage{english}

\setmainfont{DejaVu Serif}[Scale=0.92]
\setsansfont{DejaVu Sans}[Scale=0.90]
\setmonofont{DejaVu Sans Mono}[Scale=0.85]

% Кратенка за термини задржани на англиски (правилна поделба на зборови)
\newcommand{\eng}[1]{\textenglish{#1}}

% Бидејќи не се вчитани шеми за поделба на зборови за македонски,
% се олабавува порамнувањето за да се избегнат предолги редови.
\tolerance=2000
\setlength{\emergencystretch}{4em}
\hbadness=2000

%---------------------------------------------------------------
% Математика
%---------------------------------------------------------------
\usepackage{amsmath}
\usepackage{amssymb}
\usepackage{mathtools}

%---------------------------------------------------------------
% Табели, распоред, типографија
%---------------------------------------------------------------
\usepackage{geometry}
\geometry{margin=2.4cm}
\usepackage{booktabs}
\usepackage{array}
\usepackage{tabularx}
\usepackage{ragged2e}
\usepackage{caption}
\captionsetup{labelfont=bf,font=small,skip=6pt}
\usepackage{microtype}
\usepackage{enumitem}
\setlist{itemsep=2pt,topsep=3pt}
\usepackage{xcolor}
\definecolor{linknavy}{RGB}{20,60,120}

% Тип на колона со порамнување налево и автоматско прекршување
\newcolumntype{L}{>{\RaggedRight\arraybackslash}X}

%---------------------------------------------------------------
% Хиперврски (се вчитува меѓу последните)
%---------------------------------------------------------------
\usepackage{hyperref}
\hypersetup{
  colorlinks=true,
  linkcolor=black,
  citecolor=linknavy,
  urlcolor=linknavy,
  pdftitle={Дизајн на трифазен инвертор за BLDC мотор 3--5 kW},
  pdflang={mk}
}
\usepackage{url}
\urlstyle{sf}

%---------------------------------------------------------------
% Наслов
%---------------------------------------------------------------
\title{\bfseries Дизајн на трифазен инвертор со N-канални MOSFET транзистори за BLDC мотор од 3--5 kW напојуван од 24-V / 48-V батерии (опсег 20--30\,V и 40--60\,V)}

\author{Студија на ниво на компоненти со употреба на Infineon IPT015N10N5\\
OptiMOS\textsuperscript{\texttrademark}\,5 MOSFET и Texas Instruments UCC27211A-Q1 \eng{half-bridge gate driver}}

\date{}

\begin{document}
\maketitle

%===============================================================
\begin{abstract}
\noindent
Овој извештај го претставува дизајнот на енергетскиот и управувачкиот (\eng{gate-drive}) степен на трифазен напонски инвертор наменет за погон на безчеткест еднонасочен (BLDC) мотор. Системот се напојува \emph{од батерии}: пакет од 24\,V, чиј напон се менува од 20 до 30\,V во зависност од состојбата на полнење, што дава околу 3\,kW; и пакет од 48\,V, со напон од 40 до 60\,V, што дава до 5\,kW. Промената на напонот на батеријата е суштинска за дизајнот: \emph{највисокиот} напон (60\,V) ги одредува напрезите на блокирање и $dV/dt$, додека \emph{најнискиот} напон (20\,V) дава најголема струја за дадена моќност и оттука најлошиот случај за спроводните загуби. Излезниот степен користи по два Infineon IPT015N10N5 100\,V OptiMOS\textsuperscript{\texttrademark}\,5 MOSFET транзистори поврзани паралелно за секоја прекинувачка позиција (четири уреди по гранка на полумостот, дванаесет вкупно), а секоја гранка ја управува по еден Texas Instruments UCC27211A-Q1 120\,V, 3.7\,A/4.5\,A \eng{half-bridge gate driver}, напојуван од 12\,V. Целната прекинувачка фреквенција е 40\,kHz, над аудио опсегот. Во првиот дел се потврдува електричната компатибилност на избраните компоненти за целиот опсег на напони на батериите. Во вториот дел се изведуваат вредностите на пасивните помошни компоненти за управување на гејтот --- \eng{bootstrap} кондензатор, барање за \eng{bootstrap} диода, отпорници на гејтот (вклучувајќи снага и куќиште), \eng{pull-down} отпорници гејт--извор и кондензатори за раздвојување --- врз основа на објавени проектантски равенки. Во последниот дел се објаснува зошто е потребен инвертор за BLDC машина, како работат полумостот и \eng{bootstrap} напојувањето, улогата на \eng{Miller plateau} и на големите изворни/понорни струи на драјверот, како и внатрешната архитектура на UCC27211A-Q1. Секоја квантитативна тврдба и секоја проектантска равенка се поткрепени со примарен извор (каталошки лист на производителот или апликативна белешка).
\end{abstract}

%===============================================================
\section{Спецификација на дизајнот}

\begin{table}[h]
\centering
\caption{Зададени параметри на дизајнот.}
\begin{tabularx}{\textwidth}{@{}>{\RaggedRight\arraybackslash}p{3.4cm} >{\RaggedRight\arraybackslash}p{3.6cm} L@{}}
\toprule
\textbf{Параметар} & \textbf{Вредност} & \textbf{Образложение / извор} \\
\midrule
Извор на напојување & Батерии (24\,V или 48\,V пакет) & Напонот варира со состојбата на полнење на батеријата \\
Напон на еднонасочната собирница (\eng{DC-link}) & 24\,V пакет: 20--30\,V \newline 48\,V пакет: 40--60\,V & Опсег на напон на батеријата; $V_{bus,max}=60$\,V (полна 48\,V батерија), $V_{bus,min}=20$\,V (празна 24\,V батерија) \\
Излезна моќност & $\approx 3$\,kW (24\,V пакет) / $\approx 5$\,kW (48\,V пакет) & При номинален напон; моќноста на моторот е определена од напонот на собирницата \\
Прекинувачка фреквенција, $f_{SW}$ & 40\,kHz & Над границата на човечкиот слух ($\sim 20$\,kHz), па PWM тонот е нечуен \\
Прекинувачки уред & IPT015N10N5, $\times 2$ паралелно по позиција & 100\,V, 1.5\,m$\Omega$ OptiMOS\textsuperscript{\texttrademark}\,5, оптимизиран за апликации со голема струја (LEV/виљушкари/електрични алати) \cite{ref1,ref2} \\
\eng{Gate driver} & UCC27211A-Q1, $\times 1$ по фаза & 120\,V \eng{boot}, 3.7\,A/4.5\,A полумостен драјвер со вградена \eng{boot} диода \cite{ref3,ref4} \\
Напон за управување на гејтот, $V_{DD}$ & 12\,V & Во рамки на работниот опсег 8--17\,V на уредот \cite{ref3} \\
\bottomrule
\end{tabularx}
\end{table}

\noindent
Бидејќи моторот се напојува од батерија со променлив напон, струите за најлош случај мора да се проценат при \emph{најнискиот} напон на пакетот, бидејќи за дадена моќност струјата расте кога напонот опаѓа ($I = P/V$). Конзервативна граница е добиена ако се претпостави дека конверторот сè уште бара номинална моќност и при празна батерија:
\begin{itemize}
  \item Пакет 24\,V (празен, 20\,V), 3\,kW $\rightarrow I_{bus,max} = 3000/20 = 150\,\mathrm{A}$;
  \item Пакет 48\,V (празен, 40\,V), 5\,kW $\rightarrow I_{bus,max} = 5000/40 = 125\,\mathrm{A}$.
\end{itemize}
Затоа меродавната струја за најлош случај за термичкото димензионирање е $\approx 150\,\mathrm{A}$ (24\,V пакет при празна батерија). За споредба, при номинален напон струите се $5000/48 \approx 104\,\mathrm{A}$ (48\,V) и $3000/24 \approx 125\,\mathrm{A}$ (24\,V). Бидејќи во секоја прекинувачка позиција се поврзани паралелно два уреди, секој MOSFET носи приближно половина од оваа струја, т.е. до $\approx 75\,\mathrm{A}$ во најлош случај, пред да се земат предвид повисоките моментни врвови на фазната струја кај трапезоидната комутација. Ова сè уште е далеку под континуираната оценка од $2\times 250\,\mathrm{A}$ ($T_C=100\,^\circ$C) на паралелниот пар \cite{ref1}. Спроводната загуба по уред во овој најлош случај е приближно $P = I^2 R_{DS(on)} = 75^2 \times 1.5\,\mathrm{m}\Omega \approx 8.4\,\mathrm{W}$ (моментно, при спроводност), што бара соодветно ладење (в.\,\S\,\ref{sec:concl}).

%===============================================================
\section{Истражување на компонентите и проценка на компатибилноста}

\subsection{Infineon IPT015N10N5 --- клучни параметри}

IPT015N10N5 е N-канален OptiMOS\textsuperscript{\texttrademark}\,5 енергетски транзистор во \eng{TO-Leadless} (PG-HSOF-8 / „TOLL“) куќиште, кое Infineon експлицитно го опишува како „идеално за прекинување со висока фреквенција“ и наменето за погони на лесни електрични возила, виљушкари и електрични алати \cite{ref1,ref2}. Параметрите релевантни за овој дизајн се сумирани подолу; сите вредности се од каталошкиот лист рев.\,2.4, освен ако не е поинаку наведено \cite{ref1}.

\begin{table}[h]
\centering
\caption{Клучни параметри на IPT015N10N5 \cite{ref1}.}
\begin{tabularx}{\textwidth}{@{}l L L@{}}
\toprule
\textbf{Симбол} & \textbf{Параметар} & \textbf{Вредност / забелешка} \\
\midrule
$V_{DS}$ & Напон одвод--извор & 100\,V (апсолутен максимум) \\
$R_{DS(on)}$ & Отпорност во спроводна состојба при $V_{GS}=10$\,V & 1.3\,m$\Omega$ тип. / 1.5\,m$\Omega$ макс. \\
$I_D$ & Континуирана струја на одвод & 353\,A ($T_C=25\,^\circ$C); 250\,A ($T_C=100\,^\circ$C) \\
$V_{GS(th)}$ & Прагов напон на гејтот & 2.2 / 3.0 / 3.8\,V (мин./тип./макс.) \\
$V_{plateau}$ & Напон на \eng{Miller plateau} & 4.4\,V тип. (при $I_D=100$\,A) \\
$C_{iss}$ & Влезна капацитивност & 12\,nF тип. / 16\,nF макс. \\
$C_{oss}$ & Излезна капацитивност & 1.8\,nF тип. \\
$C_{rss}$ & Преносна (\eng{reverse-transfer}) капацитивност & 80\,pF тип. / 140\,pF макс. \\
$R_{G}$ & Внатрешна отпорност на гејтот & 1.4\,$\Omega$ тип. / 2.1\,$\Omega$ макс. \\
$Q_{g}$ & Вкупен полнеж на гејтот (0--10\,V) & 169\,nC тип. / 211\,nC макс. \\
$Q_{gs}$ & Полнеж гејт--извор & 53\,nC \\
$Q_{gd}$ & Полнеж гејт--одвод (\eng{Miller}) & 34\,nC тип. / 51\,nC макс. \\
$g_{fs}$ & Трансконтуктанса & 140\,S мин. / 280\,S тип. \\
$V_{SD}$ & Напон на телесната диода (\eng{body diode}) & 0.85\,V тип. / 1.2\,V макс. \\
$V_{GS,max}$ & Граничен напон гејт--извор & $\pm 20$\,V \\
\bottomrule
\end{tabularx}
\end{table}

Особина што директно го оправдува паралелното поврзување е позитивниот температурен коефициент на $R_{DS(on)}$ (околу $+0.7\,\%/^\circ$C до $+1\,\%/^\circ$C кај енергетските MOSFET транзистори): уред што спроведува поголема струја повеќе се загрева, неговата отпорност расте и струјата се враќа кон соседниот уред, така што паралелните MOSFET транзистори самостојно ги изедначуваат своите струи \cite{ref5}. Ова е стандардната причина зошто MOSFET транзисторите (за разлика од биполарните) се добро прилагодени за паралелна работа \cite{ref5}.

\subsection{Texas Instruments UCC27211A-Q1 --- клучни параметри}

UCC27211A-Q1 е AEC-Q100 квалификуван 120\,V \eng{high-side}/\eng{low-side} \eng{half-bridge gate driver}, дизајниран да управува два N-канални MOSFET транзистори со независни влезови; неговиот лебдечки \eng{high-side} степен може да работи до 120\,V, а 120\,V \eng{bootstrap} диода е интегрирана на чипот, така што не е задолжителна надворешна \eng{boot} диода \cite{ref3,ref4}. Релевантните параметри, сите од каталошкиот лист SLUSCG0C \cite{ref3}:

\begin{table}[h]
\centering
\caption{Клучни параметри на UCC27211A-Q1 \cite{ref3}.}
\begin{tabularx}{\textwidth}{@{}l L L@{}}
\toprule
\textbf{Симбол} & \textbf{Параметар} & \textbf{Вредност} \\
\midrule
$V_{DD}$ & Напојување (работно / апс.\,макс.) & 8--17\,V / 20\,V \\
$I_{O,source}$ / $I_{O,sink}$ & Врвна излезна струја & 3.7\,A / 4.5\,A \\
$V_{DDR}$ & UVLO праг на $V_{DD}$ (растечки) & 7.0\,V (хистереза 0.5\,V) \\
$V_{HBR}$ & UVLO праг на \eng{high-side} (HB--HS, растечки) & 6.7\,V (хистереза 1.1\,V) \\
$V_{HB}$ & Макс. напон на \eng{boot} напојувањето & 120\,V (HB--HS макс. 20\,V) \\
$V_{HS}$ & Опсег на јазолот за прекинување & $-(24-V_{DD})$\,V ($-12$\,V при 12\,V), до 105\,V (препорачано) \\
$V_{F}$ & Пад на напон на \eng{boot} диодата & 0.85\,V (@100\,$\mu$A) / 1.05\,V (@100\,mA) макс. \\
$R_{D}$ & Динамичка отпорност на \eng{boot} диодата & 0.55\,$\Omega$ тип. \\
$t_{PD}$ & Доцнење на простирање (\eng{propagation delay}) & 20\,ns тип. \\
$t_{R}$/$t_{F}$ & Време на пораст/опаѓање (1\,nF) & 7.2\,ns / 5.5\,ns \\
$I_{HB}$ & Струја на мирување на \eng{boot} & 0.12\,mA макс. \\
$SR_{HS}$ & Макс. \eng{slew rate} на HS & 50\,V/ns \\
$C_{HB}$ & Препорачан опсег на \eng{boot} кондензатор & 0.022--0.1\,$\mu$F (зависно од полнежот на гејтот) \\
$C_{VDD}$ & Препорачано раздвојување на $V_{DD}$ & 0.22--4.7\,$\mu$F \\
\bottomrule
\end{tabularx}
\end{table}

\subsection{Заклучок за компатибилноста}

\begin{table}[h]
\centering
\caption{Проверка на компатибилноста на компонентите.}
\begin{tabularx}{\textwidth}{@{}L L L >{\RaggedRight\arraybackslash}p{2.1cm}@{}}
\toprule
\textbf{Проверка} & \textbf{Барање} & \textbf{Овој дизајн} & \textbf{Исход} \\
\midrule
Напон на блокирање & $V_{DS} \gg V_{bus,max}$ вкл.\,\eng{ringing} & 100\,V наспроти 60\,V (полна 48\,V батерија); $\approx 1.67\times$ DC маргина & \textbf{Поминува} --- $\approx 40$\,V резерва за надскок при исклучување \cite{ref1} \\
Напон на \eng{high-side} напојувањето & $V_{HB} > V_{bus,max}$ & 120\,V \eng{boot}, 105\,V препорачано HS, наспроти 60\,V & \textbf{Поминува} \cite{ref3} \\
Ниво на управување на гејтот & $V_{GS,max} > V_{DD} \ge$ полна спроводливост & $\pm 20$\,V граница, управуван на 12\,V & \textbf{Поминува} \cite{ref1,ref3} \\
UVLO & $V_{DD} >$ UVLO & 12\,V наспроти 7\,V ($V_{DD}$), 6.7\,V (HB) & \textbf{Поминува} \cite{ref3} \\
Негативно \eng{ringing} на јазолот & HS поднесува потсек & $-12$\,V преодна способност при $V_{DD}=12$\,V & \textbf{Поминува} \cite{ref3} \\
Струен капацитет & $I_D \gg$ струја на гранката & $2\times 250$\,A (100\,$^\circ$C) наспроти $\approx 75$\,A (најлош случај при $V_{bus,min}=20$\,V) & \textbf{Поминува} \cite{ref1} \\
Сила на управување наспроти полнеж & Доволна струја низ \eng{Miller plateau} за \emph{паралелниот} полнеж & в.\,\S\,\ref{sec:gateR} & \textbf{Поминува} \\
\bottomrule
\end{tabularx}
\end{table}

\noindent
Спрегата е електрично исправна за целиот опсег на напони на батериите. Со највисокиот напон на собирницата (60\,V при полна 48\,V батерија), 100\,V класата дава DC маргина од $\approx 1.67\times$ и околу 40\,V резерва за надскокот на $V_{DS}$ при исклучување предизвикан од $L\cdot di/dt$; ова е поудобно отколку би изгледало, но прави управувањето на брзината на исклучување (преку асиметричната мрежа на отпорници на гејтот во \S\,\ref{sec:asym}) уште поважно одошто на фиксна 48\,V собирница. Единствената точка што бара квантитативна проверка --- бидејќи задачата паралелно поврза два MOSFET транзистори со висок $Q_g$ по прекинувач --- е дали еден канал на UCC27211A-Q1 може доволно брзо да го наполни и испразни \emph{удвоениот} полнеж на гејтот. Ова е разгледано во \S\,\ref{sec:gateR} и во димензионирањето на \eng{bootstrap}-от во \S\,\ref{sec:cboot}, и е причината зошто стандардниот 0.1\,$\mu$F \eng{bootstrap} кондензатор препорачан во каталошкиот лист мора да се преиспита за оваа оптоварувачка состојба.

%===============================================================
\section{Дизајн и пресметки на пасивните компоненти}
\label{sec:passive}

Низ целиот текст, полнежот на гејтот за една \emph{прекинувачка позиција} е двојно поголем од оној на еден уред, бидејќи се поврзани паралелно два MOSFET транзистори:
\begin{equation}
Q_{g,\text{поз}} = 2 \times Q_{g} = 2 \times 211\,\mathrm{nC} = 422\,\mathrm{nC}\ \text{(најлош случај, до 10\,V)} \quad \cite{ref1}
\end{equation}
Бидејќи гејтот се управува до 12\,V наместо до 10\,V (при кои е специфициран $Q_g$), и бидејќи кривата на полнежот е речиси рамна над $\sim 7$\,V (па дополнителниот полнеж од 10 до 12\,V е претежно линеарниот член $C_{GS}$), се додава резерва од $\approx 10\,\%$. Дополнително, треба да се забележи дека $Q_g$ е специфициран при напон одвод--извор од 50\,V \cite{ref1}; при најлошиот напон на собирницата од 60\,V, \eng{Miller} полнежот $Q_{gd}$ (кој зависи од нишањето на напонот на одводот) е маргинално поголем (за неколку nC по уред бидејќи $C_{rss}$ нагло опаѓа со напонот), но ова исто така е опфатено со истата резерва. Резултирачката проектна вредност е:
\begin{equation}
Q_{g,\text{поз,проект}} \approx 465\,\mathrm{nC}\ \text{по прекинувачка позиција.}
\end{equation}

\subsection{Bootstrap кондензатор}
\label{sec:cboot}

\paragraph{Метод.} \eng{Bootstrap} кондензаторот $C_{BOOT}$ го обезбедува целиот полнеж потребен за вклучување на \eng{high-side} уредот и за напојување на струјата на мирување на \eng{high-side} степенот за времето на спроводност на \eng{high-side}-от, додека опаѓа за не повеќе од дозволената вредност $\Delta V_{BS}$ \cite{ref6,ref7}. Управувачките равенки се:
\begin{align}
\Delta V_{BS} &\le V_{DD} - V_{F} - V_{GS,\min} - V_{DS,on} \label{eq:dvbs}\\
Q_{TOTAL} &= Q_{g,\text{поз}} + \frac{I_{QBS} + I_{LK}}{f_{SW}} + Q_{LS}, \qquad C_{BOOT} \ge \frac{Q_{TOTAL}}{\Delta V_{BS}} \label{eq:cboot}
\end{align}
со дополнителна резерва од $\sim 20\,\%$ препорачана за толеранции и опаѓање на капацитивноста кога диодата е интегрирана \cite{ref9}.

\paragraph{Бројки.}
\begin{itemize}
  \item Пад на \eng{boot} диодата при релевантната струја на полнење: $V_F \approx 0.9$\,V \cite{ref3}.
  \item Се бара гејтот на \eng{high-side}-от да остане целосно спроводен, $V_{GS,\min} \approx 10$\,V; \eng{low-side} $V_{DS,on}\approx 0.1$\,V. Тогаш \emph{дозволеното} опаѓање според \eqref{eq:dvbs} е $\Delta V_{BS} \le 12 - 0.9 - 10 - 0.1 = 1.0\,\mathrm{V}$ \cite{ref6}. Се избира конзервативна проектна цел од \textbf{$\Delta V_{BS}=0.5$\,V}, што го задржува лебдечкото напојување на $\approx 10.6$\,V --- далеку над 6.7\,V UVLO прагот на HB \cite{ref3}.
  \item Полнеж од мирување/протекување во текот на најлошото време на спроводност на \eng{high-side}. Со $I_{QBS}=0.12$\,mA \cite{ref3}, MOSFET $I_{GSS}=2\times 100$\,nA \cite{ref1} и висок \eng{duty cycle} $D_{max}\approx 0.95$, па $t_{HS}=0.95/40\,\mathrm{kHz}=23.75\,\mu\mathrm{s}$:
  \begin{equation}
  Q_{leak} \approx (0.12\,\mathrm{mA})(23.75\,\mu\mathrm{s}) \approx 2.9\,\mathrm{nC} \ll Q_{g,\text{поз}}.
  \end{equation}
\end{itemize}
Полнежот на гејтот целосно доминира, па:
\begin{equation}
Q_{TOTAL} \approx 465\,\mathrm{nC} + 3\,\mathrm{nC} \approx 468\,\mathrm{nC}
\end{equation}
\begin{equation}
\boxed{\,C_{BOOT} \ge \dfrac{468\,\mathrm{nC}}{0.5\,\mathrm{V}} \approx 0.94\,\mu\mathrm{F}\,}
\end{equation}
Со примена на препорачаната резерва и со отчитување на опаѓањето на капацитивноста при DC-преднапон кај керамичките кондензатори од класа\,II, избраната вредност е:
\begin{quote}
\textbf{$C_{BOOT} = 2.2\,\mu$F, X7R, $\ge 25$\,V, керамички со ниска ESR}, поставен директно меѓу пиновите HB--HS \cite{ref3,ref9}.
\end{quote}

\paragraph{Забелешка за опсегот 0.022--0.1\,$\mu$F од каталошкиот лист.} Тој опсег е точен за \emph{еден} MOSFET со $\sim 30$--90\,nC полнеж на гејтот \cite{ref3}; овде паралелниот пар бара $\approx 5\times$ повеќе полнеж, што е токму причината зошто каталошкиот лист наведува дека вредноста „зависи од полнежот на гејтот на \eng{high-side} MOSFET-от“ \cite{ref3}. Со употреба на 0.1\,$\mu$F би се добило опаѓање од $468\,\mathrm{nC}/0.1\,\mu\mathrm{F} \approx 4.7$\,V по циклус --- доволно за недоволно управување или дури активирање на \eng{high-side} UVLO. Затоа поголемата вредност од 2.2\,$\mu$F не е опционална за оваа оптоварувачка состојба.

\paragraph{Проверка на полнењето.} Кондензаторот го надополнува само $\approx 468$\,nC отстранет во секој циклус, не целиот полнеж. При најлош \eng{duty cycle}, времето на спроводност на \eng{low-side}-от е $\approx 1.25\,\mu$s; потребната просечна струја на полнење е $468\,\mathrm{nC}/1.25\,\mu\mathrm{s} \approx 0.37$\,A, додека интегрираната 0.55\,$\Omega$ \eng{boot} диода може почетно да испорача $(12-0.9)/0.55 \approx 20$\,A \cite{ref3}. Затоа полнењето е удобно при 40\,kHz. При многу мала брзина на моторот со \eng{duty cycle} близу 100\,\% на \eng{high-side}-от, прозорецот за надополнување се намалува и може да биде потребен периодичен \eng{low-side} освежувачки импулс --- позната ограниченост на сите \eng{bootstrap} напојувања \cite{ref6,ref8}.

\subsection{Bootstrap диода}

UCC27211A-Q1 ја \emph{интегрира} \eng{bootstrap} диодата (анода кон $V_{DD}$, катода кон HB); таа е оценета на 120\,V, има типична динамичка отпорност од 0.55\,$\Omega$, пад на напон од $\approx 0.9$--1.05\,V при 100\,mA и време на исклучување од 20\,ns \cite{ref3}. За овој дизајн внатрешната диода е соодветна: \S\,\ref{sec:cboot} покажува дека и струјата на полнење и опаѓањето се удобно во рамки на нејзината способност, а нејзината 120\,V оценка ја надминува 48\,V собирницата со голема маргина \cite{ref3}.

\emph{Опционална надворешна диода паралелно} со внатрешната може да се додаде ако проектантот сака дополнително да го намали падот на напон при полнење и опаѓањето по циклус. Ако се користи, мора да задоволи \cite{ref6,ref8}: обратен напон $\ge V_{bus} + V_{DD}$ маргина $\rightarrow$ се избира \textbf{100\,V} компонента за совпаѓање со напонската класа на MOSFET-от; брзо обратно закрепнување (\eng{reverse recovery}) ($t_{rr}\le 50$\,ns) или \eng{Schottky}; и струја на спроводност $\ge$ врвната струја на полнење ($\ge 1$\,A е доволно). Бидејќи функцијата веќе е обезбедена на чипот, оваа компонента е надградба, а не неопходност.

\subsection{Отпорници на гејтот и проверка на врвната струја / Miller}
\label{sec:gateR}

\paragraph{Барање за пригушување (Balogh).} Оптималниот надворешен отпорник на гејтот за критично пригушување на резонанцата во колото на гејтот, формирана од индуктивноста на изворот $L_S$ и $C_{iss}$, е \cite{ref5}:
\begin{equation}
R_{GATE,OPT} = 2\sqrt{\frac{L_S}{C_{ISS}}} - (R_{DRV} + R_{G,I})
\end{equation}
За куќиштето TOLL со ниска индуктивност ($L_S$ од редот на 1\,nH) и $C_{iss}=12$\,nF, $2\sqrt{L_S/C_{ISS}} \approx 0.58\,\Omega$, додека излезната импеданса на драјверот плус внатрешната отпорност на гејтот ($R_{DRV}+R_{G,I}$) веќе ја надминуваат. Резултатот е негативен, што значи дека колото е \textbf{веќе прекумерно пригушено само од импедансата на драјверот и внатрешната отпорност на гејтот од 1.4\,$\Omega$} \cite{ref5,ref1}. Затоа надворешниот отпорник на гејтот се избира не за пригушување, туку за: (i) да ја држи споделената излезна струја на драјверот во рамки на оценката, (ii) да го наштелува $dV/dt$ заради EMI и (iii) да ги избалансира двата паралелни уреди.

\paragraph{Поединечни отпорници.} Најдобрата практика за паралелни MOSFET транзистори е \emph{поединечен} отпорник на гејтот кај секој гејт (по два по прекинувачка позиција) наместо еден заеднички отпорник, така што секоја тенденција двата уреди да осцилираат еден против друг се пригушува и нивното прекинување се балансира. Се избира стандардна вредност:
\begin{quote}
\textbf{$R_{G,ext} = 2.2\,\Omega$ (по еден на секој MOSFET, E24)},
\end{quote}
што е блиску до 1.8\,$\Omega$ при кои се карактеризирани времињата на прекинување на IPT015N10N5 \cite{ref1}.

\paragraph{Проверка на врвната струја / $dV/dt$.} \eng{Miller plateau} е на $V_{plateau}=4.4$\,V \cite{ref1}. Со гејтот задржан на платото, мрежата би \emph{барала}, по уред,
\begin{equation}
I_{G} = \frac{V_{DD} - V_{plateau}}{R_{G,ext} + R_{G,I}} = \frac{12 - 4.4}{2.2 + 1.4} \approx 2.1\,\mathrm{A},
\end{equation}
т.е.\,$\approx 4.2$\,A за двата уреди заедно. Ова ја надминува границата на изворната струја на драјверот од 3.7\,A \cite{ref3}, па драјверот ја ограничува струјата на 3.7\,A и напонот на јазолот опаѓа --- драјверот, а не отпорникот, ја одредува брзината. Времето за испорака на комбинираниот \eng{Miller} полнеж ($2\times 51\,\mathrm{nC}=102$\,nC макс.) при 3.7\,A е:
\begin{equation}
t_{Miller} = \frac{Q_{gd,\text{поз}}}{I_{O,source}} = \frac{102\,\mathrm{nC}}{3.7\,\mathrm{A}} \approx 28\,\mathrm{ns},
\end{equation}
што дава, на најлошиот случај на собирницата од 60\,V (полна 48\,V батерија),
\begin{equation}
\frac{dV_{DS}}{dt} \approx \frac{60\,\mathrm{V}}{28\,\mathrm{ns}} \approx 2.1\,\mathrm{V/ns},
\end{equation}
што е умерено (добро за спроводено EMI) и далеку под границата на \eng{slew rate} од 50\,V/ns на пинот HS \cite{ref3}.

Ова го потврдува клучното прашање за компатибилноста: \textbf{дури и со удвоениот полнеж на гејтот од паралелното поврзување, еден канал на UCC27211A-Q1 го прекинува парот за под 30\,ns} --- токму за што е наменета оценката 3.7\,A/4.5\,A на уредот, имено „управување на големи енергетски MOSFET транзистори со минимизирани прекинувачки загуби при преминот низ \eng{Miller plateau}-то“ \cite{ref3}. Сепак, единствената вредност од 2.2\,$\Omega$ изведена овде е само \emph{симетрична основа}: \S\,\ref{sec:asym} ја префинува во асиметрична мрежа за вклучување/исклучување, која е конфигурацијата што всушност е усвоена.

%---------------------------------------------------------------
\subsection{Асиметрична (поделена) мрежа на отпорници на гејтот}
\label{sec:asym}

Единствениот отпорник од \S\,\ref{sec:gateR} принудува една вредност да служи за двата фронта на прекинување, што е компромис. Овој потоддел го заменува со \textbf{асиметрична} мрежа што ги поставува брзините на вклучување и исклучување независно --- побавно, помеко вклучување и брзо исклучување со ниска импеданса --- што е конфигурацијата усвоена во финалниот дизајн.

\paragraph{Топологија.} По MOSFET, секогаш присутниот сериски отпорник \textbf{$R_{on}$} ја носи струјата на вклучување (изворна), а гранка од \textbf{$R_{off}$ во серија со OFF диода $D_{off}$} се поставува паралелно со него. $D_{off}$ е ориентирана со анода кон гејтот и катода кон излезот на драјверот, така што спроведува \emph{само кога гејтот се влече надолу} (исклучување / $dV/dt$ стегач). Следствено:
\begin{equation}
R_{g(\text{on})} = R_{on}, \qquad R_{g(\text{off})} = R_{on} \parallel (R_{off} + R_{D})
\end{equation}
При вклучување $D_{off}$ е обратно поларизирана и гејтот гледа само $R_{on}$; при исклучување двете патеки спроведуваат и гејтот ја гледа пониската паралелна комбинација \cite{ref5,ref12}. (Со една OFF диода, $R_{on}$ останува во патеката на исклучување; пониската гранка $R_{off}$ едноставно го шунтира, поради што $R_{g(off)}$ е паралелна комбинација, а не само $R_{off}$.)

\paragraph{Проектна насока --- бавно вклучување, брзо исклучување.} За тврдо прекинуван полумост, претпочитаната асиметрија е \emph{побавно вклучување, побрзо исклучување} \cite{ref12,ref13,ref14}. Побрзото вклучување го зголемува $dI/dt$ на струјата на одводот, што ја зголемува струјата на обратно закрепнување (\eng{reverse recovery}) на \emph{телесната \eng{freewheeling} диода на комплементарниот уред}; телесната диода на IPT015N10N5 складира $Q_{rr}\approx 316$\,nC ($t_{rr}\approx 103$\,ns) \cite{ref1}, па намерно омекнато вклучување го ограничува тој врв на закрепнување, придружната загуба и EMI што таа го генерира \cite{ref12,ref13}. Спротивно на тоа, исклучувањето има корист од тоа да биде брзо и со ниска импеданса: ги намалува загубите при исклучување и, што е клучно, го држи гејтот на \emph{исклучениот} уред цврсто кон масата, така што \eng{Miller} струјата инјектирана низ $C_{GD}$ кога комплементарниот прекинувач го менува јазолот не може да го подигне $V_{GS}$ до прагот и да предизвика само-(лажно) вклучување \cite{ref13,ref14}. 100\,V уредот на собирница до 60\,V има доволна напонска резерва ($\approx 40$\,V) да го апсорбира малку повисокиот $di/dt$ надскок при исклучување што го произведува брзото исклучување; токму затоа што таа резерва е помала отколку на фиксна 48\,V собирница, можноста независно да се контролира фронтот на исклучување е уште поценета. Асиметријата ја засилува и сопствената нерамнотежа 4.5\,A-понор / 3.7\,A-извор на драјверот \cite{ref3}.

\paragraph{Избор на вредности.}
\begin{itemize}
  \item \textbf{$R_{on} = 3.3\,\Omega$ (по MOSFET).} Зголемен над компромисот од 2.2\,$\Omega$ од \S\,\ref{sec:gateR} за да се омекне фронтот на вклучување. Резултирачката струја на вклучување по уред на платото е $I=(12-4.4)/(R_{on}+R_{G,I})=7.6/(3.3+1.4)=1.62$\,A, т.е.\,$\approx 3.2$\,A за парот --- сега поставена од отпорникот, а не од границата на драјверот од 3.7\,A, што дава контролиран $dV/dt$ од $60\,\mathrm{V}/(102\,\mathrm{nC}/3.2\,\mathrm{A})\approx 1.9$\,V/ns во најлош случај (60\,V собирница) \cite{ref1,ref3}.
  \item \textbf{$R_{off} = 1.0\,\Omega$ (по MOSFET).} Со динамичка отпорност на OFF диодата $R_{D}\approx 0.1$--0.15\,$\Omega$ (проценета од спроводната карактеристика на PMEG4030ER, $V_F\approx 0.46$--0.54\,V во релевантниот опсег на струја \cite{ref12}):
  \begin{equation}
  R_{g(\text{off})} = R_{on} \parallel (R_{off}+R_D) = 3.3 \parallel 1.15 \approx 0.85\,\Omega.
  \end{equation}
  Ова дава ефективен однос на отпорност вклучување:исклучување од околу \textbf{3.9\,:\,1}.
\end{itemize}

\paragraph{Зошто $R_{off}$ не се намалува понатаму.} \emph{Внатрешната} отпорност на гејтот $R_{G,I}=1.4\,\Omega$ е во серија со чипот и не може да се заобиколи \cite{ref1}. Затоа вкупната отпорност на колото на гејтот при исклучување е $R_{G,I}+R_{g(\text{off})}=1.4+0.85=2.25\,\Omega$, и дури и со $R_{off}=0$ не би можела да падне под $\approx 1.4\,\Omega$. Намалувањето на $R_{off}$ под $\sim 1\,\Omega$ затоа дава сè помали придобивки, па 1.0\,$\Omega$ е добро поставен избор. При оваа вредност, врвната струја на исклучување по уред на платото е $4.4/2.25 \approx 1.96$\,A ($\approx 3.9$\,A за парот), блиску до границата на понорната струја на драјверот од 4.5\,A, што дава транзиција на исклучување од $\approx 26$\,ns наспроти $\approx 32$\,ns при вклучување --- посакуваната асиметрија \cite{ref3}.

\paragraph{Избор на OFF диода.} $D_{off}$ мора: (i) да го блокира обратниот напон што се појавува на $R_{on}$ при вклучување ($\approx I_{on,peak}\cdot R_{on} \approx 1.62\,\mathrm{A}\times 3.3\,\Omega \approx 5.3$\,V, плус \eng{ringing}); (ii) да ја носи струјата на импулсот при исклучување ($\approx 2$\,A врв по гејт); (iii) да закрепнува моментално, така што не додава полнеж на следното вклучување; и (iv) да има ниска динамичка отпорност, така што $R_{off}$ доминира во патеката на исклучување. \eng{Schottky} диода ги задоволува сите четири. Избраната компонента е:
\begin{quote}
\textbf{$D_{off}$ = Nexperia PMEG4030ER} --- 40\,V, 3\,A \eng{Schottky} исправувачка диода со низок $V_{F}$ ($V_{F}\approx 0.46$--0.54\,V, $I_{FSM}=50$\,A, SOD-123W) \cite{ref12}.
\end{quote}
Нејзината оценка од 40\,V дава $>7\times$ маргина над $\sim 5$\,V обратен напрег; нејзината оценка од 3\,A просечно / 50\,A ударно ја надминува струјата на импулсот при исклучување од $\sim 2$\,A; нејзиниот низок $V_F$ и мала спојна капацитивност ја држат патеката на исклучување чиста; а како \eng{Schottky} има занемарливо обратно закрепнување \cite{ref12}. Широко достапна алтернатива е 1N5819 (40\,V, 1\,A \eng{Schottky}), со прифаќање на нејзиното поголемо куќиште и капацитивност.

\paragraph{Снага и куќиште на отпорниците.} Просечната снага дисипирана во отпорникот на гејтот е дел од вкупната загуба за управување на гејтот, $P_G = Q_g\,V_{DD}\,f_{SW}$ \cite{ref5}. По уред, со $Q_g \approx 232$\,nC при 12\,V, $P_G \approx 232\,\mathrm{nC}\times 12\,\mathrm{V}\times 40\,\mathrm{kHz} \approx 0.11$\,W, поделена подеднакво меѓу вклучувањето и исклучувањето. Уделот што паѓа на надворешниот отпорник е неговата фракција од вкупната отпорност во колото на гејтот \cite{ref5}; со проценета излезна отпорност на драјверот $R_{HI}\approx 3.2\,\Omega$:
\begin{align}
P_{R_{on}} &\approx \frac{R_{on}}{R_{HI}+R_{on}+R_{G,I}}\cdot \tfrac{1}{2}Q_g V_{DD} f_{SW} = \frac{3.3}{7.9}\times 0.0557\,\mathrm{W} \approx 26\,\mathrm{mW}, \\
P_{R_{off}} &\approx 6\,\mathrm{mW}\ \text{(само при исклучување; останатото го носат $R_{HI}$ и $R_{G,I}$).}
\end{align}
Двете се далеку под $1/16$\,W. Сепак, отпорникот гледа кратки струјни врвови ($\sim 1.6$--2\,A), па изборот на куќиште се води од енергијата по импулс и од маргината, а не само од просечната снага: енергијата по транзиција во $R_{on}$ е $\approx 0.42\times \tfrac{1}{2}Q_g V_{DD} \approx 0.6\,\mu$J, занемарливо во споредба со импулсната издржливост на кој било чип-отпорник. Се избира:
\begin{quote}
\textbf{$R_{on}$, $R_{off}$: куќиште 0805, дебелослоен (\eng{thick-film}), $1/8$\,W (125\,mW)}, што дава $>4\times$ маргина на просечната снага и солидна импулсна резерва; работниот напон ($\sim 12$\,V) е далеку под напонската оценка на 0805 ($\sim 150$\,V) \cite{ref5}.
\end{quote}
Куќиштето 0603 ($1/10$\,W) исто така е електрично доволно; 0805 е претпочитано заради подобро ширење на топлината, импулсна робусност и леснотија на монтажа. (\eng{Pull-down} отпорникот $R_{GS}$ е димензиониран во \S\,\ref{sec:pulldown}.)

\begin{table}[h]
\centering
\caption{Споредба: единствен отпорник наспроти асиметрична мрежа.}
\begin{tabularx}{\textwidth}{@{}L L L@{}}
\toprule
\textbf{Аспект} & \textbf{Единствен $R_{G}=2.2\,\Omega$ (\S\,\ref{sec:gateR})} & \textbf{Асиметрична $R_{on}=3.3\,\Omega$ / $R_{off}=1.0\,\Omega$ + $D_{off}$} \\
\midrule
Фронт при вклучување & Среден --- принуден компромис & Помек ($R_{g(on)}=3.3\,\Omega$); $dV/dt\approx 1.5$\,V/ns \cite{ref1,ref3} \\
Фронт при исклучување & Ист како при вклучување (не може да се разликува) & Брз ($R_{g(off)}\approx 0.85\,\Omega$ надв.); близу полн понор од 4.5\,A \cite{ref3} \\
Обратно закрепнување на комплем.\,диода & Врзано за изборот при исклучување & Независно укротено од бавното вклучување \cite{ref12,ref13} \\
Импеданса на стегачот во исклучена состојба & $R_{G,I}+2.2=3.6\,\Omega$ & $R_{G,I}+0.85=2.25\,\Omega \rightarrow$ повисок праг на $dV/dt$ за самовклучување \cite{ref13,ref14} \\
Независна оптимизација на фронтовите & Невозможна & Да --- целта на мрежата \cite{ref12} \\
Загуба при исклучување/прекинување & Повисока ако $R_G$ се зголеми за меко вклучување & Пониска (задржано брзо исклучување) \\
Дополнителни компоненти по MOSFET & --- & $+1$ отпорник, $+1$ \eng{Schottky} \\
\bottomrule
\end{tabularx}
\end{table}

\paragraph{Зошто асиметричниот пристап е подобар и претпочитан во овој дизајн.} Со единствен отпорник, проектантот се соочува со неизбежен конфликт \cite{ref12,ref13}: \emph{мала} вредност дава брзо исклучување (ниска загуба, силен стегач), но и брзо вклучување што го влошува обратното закрепнување на комплементарната телесна диода, $dV/dt$ и EMI; \emph{голема} вредност го омекнува вклучувањето, но истовремено го забавува исклучувањето, зголемувајќи ја загубата при исклучување и слабеејќи го $dV/dt$ стегачот што спречува \eng{shoot-through}. Асиметричната мрежа го раствора овој конфликт со раздвојување на двата фронта \cite{ref12}. Тоа раздвојување е особено вредно во \emph{оваа} апликација, бидејќи секоја од трите гранки е тврдо прекинуван полумост во кој секое вклучување комутира \eng{freewheeling} телесна диода што носи струја од редот на 100\,A; бидејќи два паралелни MOSFET транзистори со висок $Q_g$ по прекинувач ја удвојуваат и \eng{Miller} струјата и цената на секоја вкрстена спроводност; и бидејќи напонската резерва 100\,V/48\,V му дозволува на дизајнот да искористи брзо исклучување без грижа за надскок. Цената --- дванаесет дополнителни отпорници и дванаесет \eng{Schottky} диоди за целиот инвертор --- е занемарлива наспроти добивките во ефикасноста, EMI однесувањето и отпорноста на паразитско вклучување \cite{ref12,ref13,ref14}. Поради овие причини, асиметричната мрежа ја заменува основата со единствен отпорник од \S\,\ref{sec:gateR} и е конфигурацијата пренесена во финалниот список на материјали.

%---------------------------------------------------------------
\subsection{Pull-down отпорници гејт--извор}
\label{sec:pulldown}

Отпорник од гејт кон извор го држи секој MOSFET дефинирано исклучен пред излезите на драјверот да станат активни при вклучување на напојувањето и обезбедува $dV/dt$ заштита од \eng{Miller}-индуцирано вклучување \cite{ref5}. Вредноста на \eng{pull-down}-от следи од релацијата за $dV/dt$ имунитет \cite{ref5}:
\begin{equation}
\left.\frac{dv}{dt}\right|_{LIMIT} = \frac{V_{TH} - 0.007\,(T_J - 25)}{R_{G,I}\,C_{GD}}
\end{equation}
Во нормална работа, ниската излезна импеданса на драјверот и понорот од 4.5\,A доминираат во оваа заштита; главната задача на \eng{pull-down}-от е состојбата при вклучување на напојувањето. Стандардна:
\begin{quote}
\textbf{$R_{GS} = 10\,$k$\Omega$ (по еден на секој MOSFET)}
\end{quote}
е доволно голема за да не го оптоварува драјверот и доволно мала за да го држи гејтот надолу при вклучување. Проценката со капацитивен делител на индуцираното вклучување ја потврдува маргината: кога спротивниот прекинувач наметнува најлош напон од 60\,V на исклучениот уред, најлошиот пораст на гејтот низ делителот $C_{rss}/C_{iss}=80\,\mathrm{pF}/12\,\mathrm{nF}\approx 0.0067$ би бил $\approx 0.40$\,V дури и со \emph{лебдечки} гејт --- сè уште далеку под прагот од 2.2\,V \cite{ref1} --- а активниот понор од 4.5\,A ја држи вистинската вредност уште пониска \cite{ref3}.

\paragraph{Снага и куќиште на $R_{GS}$.} Кога гејтот е управуван на 12\,V, $R_{GS}$ непрекинато дисипира $P = V_{GS}^2/R_{GS} = 12^2/10\,\mathrm{k}\Omega = 14.4$\,mW (просечно помалку, бидејќи гејтот не е секогаш висок). Затоа е доволно куќиште \textbf{0603, $1/10$\,W (100\,mW)} ($>6\times$ маргина); за единственост на спецификацијата може да се користи и 0805, $1/8$\,W заедно со отпорниците на гејтот.

%---------------------------------------------------------------
\subsection{Кондензатори за раздвојување (decoupling)}

\paragraph{$V_{DD}$ (\eng{low-side}) раздвојување.} Секогаш кога LO излезот испорачува струја на гејтот, еднаков струен импулс се влече низ $V_{DD}$, па керамички кондензатор со ниска ESR мора да стои директно меѓу $V_{DD}$--$V_{SS}$ \cite{ref3}. Димензионирано за толерантно бранување $\Delta V$ со истиот аргумент за полнеж како кај \eng{bootstrap}-от:
\begin{equation}
C_{VDD} \ge \frac{Q_{g,\text{поз}}}{\Delta V} = \frac{465\,\mathrm{nC}}{0.5\,\mathrm{V}} \approx 0.93\,\mu\mathrm{F}.
\end{equation}
TI препорачува 0.22--4.7\,$\mu$F и горниот крај за ниски температури и висок полнеж на гејтот \cite{ref3}. Избрано:
\begin{quote}
\textbf{$C_{VDD} = 4.7\,\mu$F X7R + 0.1\,$\mu$F керамички, по драјвер}, со посебен кондензатор за секој од трите драјвери \cite{ref3}.
\end{quote}

\paragraph{HB (\eng{high-side}) раздвојување.} Оваа улога ја извршува самиот \eng{bootstrap} кондензатор (\S\,\ref{sec:cboot}), поставен директно меѓу пиновите HB--HS; посебната препорака од TI за 0.022--0.1\,$\mu$F е \emph{минималното} високофреквентно раздвојување и овде е надмината од пресметаната вредност од 2.2\,$\mu$F \cite{ref3}.

\paragraph{Бук на 12\,V собирницата.} Споделен бук резервоар (на пр.\,22--47\,$\mu$F електролитски или керамички) на помошната 12\,V собирница ги напојува трите драјвери и ги апсорбира збирните струјни импулси за управување на гејтот; TI советува посебен $V_{DD}$ кондензатор за раздвојување по драјвер во системи со повеќе драјвери \cite{ref3}.

\paragraph{Раздвојување на DC-link.} Бидејќи изворот е батерија поврзана преку сноп кабли, импедансата на изворот (внатрешната отпорност на батеријата плус индуктивноста на снопот) е значителна; затоа локалното \eng{DC-link} раздвојување е уште поважно отколку кај крута лабораториска собирница. Енергетскиот степен носи бранувачка струја од редот на 100\,A (до $\approx 150$\,A во најлош случај при празна 24\,V батерија) и мора да се раздвои со \eng{DC-link} банка со ниска ESR/ESL (фолиски и/или керамички) близу полумостот, димензионирана за RMS бранувачката струја и за да го држи надскокот на собирницата во рамки на оценката од 100\,V при комутација. Бидејќи највисокиот напон на батеријата е 60\,V, кондензаторите на собирницата треба да имаат напонска оценка $\ge 80$\,V (фолиски) или $\ge 100$\,V (керамички, со отчитување на опаѓањето при DC-преднапон). Прецизното димензионирање зависи од индуктивноста на снопот и е задача на ниво на плоча, а не на ниво на вредност на компонента.

%---------------------------------------------------------------
\subsection{Список на компоненти за gate-drive (по фаза / по плоча)}

\begin{table}[h]
\centering
\caption{Список на пасивни помошни компоненти за управување на гејтот.}
\begin{tabularx}{\textwidth}{@{}l L L c c@{}}
\toprule
\textbf{Ознака} & \textbf{Функција} & \textbf{Вредност} & \textbf{По гранка} & \textbf{3 гранки} \\
\midrule
$C_{BOOT}$ & \eng{Bootstrap} кондензатор (HB--HS) & 2.2\,$\mu$F X7R $\ge 25$\,V & 1 & 3 \\
$C_{VDD}$ & Раздвојување на $V_{DD}$ & 4.7\,$\mu$F X7R + 0.1\,$\mu$F & 1 сет & 3 сета \\
$R_{on}$ & Отпорник за вклучување (по MOSFET) & 3.3\,$\Omega$, 0805, $1/8$\,W & 4 (2 HS + 2 LS) & 12 \\
$R_{off}$ & Отпорник за исклучување (по MOSFET) & 1.0\,$\Omega$, 0805, $1/8$\,W & 4 & 12 \\
$D_{off}$ & OFF \eng{Schottky} диода & Nexperia PMEG4030ER (40\,V, 3\,A) & 4 & 12 \\
$R_{GS}$ & \eng{Pull-down} гејт--извор & 10\,k$\Omega$, 0603, $1/10$\,W & 4 & 12 \\
$D_{BOOT}$ (опц.) & Надворешна паралелна \eng{boot} диода & 100\,V брза/\eng{Schottky} $\ge 1$\,A & 0--1 & 0--3 \\
$C_{BULK}$ & Бук на 12\,V собирница & 22--47\,$\mu$F & --- & 1 (споделен) \\
\bottomrule
\end{tabularx}
\end{table}

\noindent
Плус активните уреди: \textbf{12 $\times$ IPT015N10N5} (4 по гранка) и \textbf{3 $\times$ UCC27211A-Q1}.

%===============================================================
\section{Позадина и теорија на работата}

\subsection{Зошто се користи инвертор за погон на BLDC мотор}

Безчеткестиот еднонасочен мотор е, и покрај своето име, синхрона машина за наизменична струја: има ротор со траен магнет и трифазна статорска намотка, и нема механички комутатор и четки што ги користи конвенционалниот еднонасочен мотор за префрлање на струјата меѓу намотките додека роторот се врти \cite{ref10,ref11}. Затоа функцијата на комутација што некогаш ја вршеле четките механички, мора да се изврши електронски. Инверторот е колото што го прави тоа: го претвора еднонасочниот напон на собирницата со фиксен поларитет (овде обезбеден од батерија) во низа контролирани, наизменични фазни струи потребни за да се одржи статорското магнетно поле да се врти пред роторот и со тоа да се произведува континуиран вртежен момент \cite{ref10,ref11}.

Во трапезоидната (шестчекорна, 120$^\circ$) контрола, најчестата BLDC шема, шесте прекинувачи на инверторот заземаат шест комбинации на вклучено/исклучено, произведувајќи шест дискретни ориентации на статорското поле по електрична револуција; прекинувачите се чекорат низ оваа низа синхроно со роторот, така што вртежниот момент е секогаш позитивен \cite{ref10}. Позицијата на роторот --- потребна за да се знае \emph{кога} да се чекори --- се добива или од Холови сензори поставени на 120$^\circ$, или, без сензори, од \eng{back-EMF} (контраелектромоторна сила) на моментно неенергизираната фаза \cite{ref10,ref11}. Во секој момент една фаза се управува позитивно, една негативно, а една се остава отворена, совпаѓајќи ја правоаголната струја со трапезоидната \eng{back-EMF} на моторот за ефикасно производство на вртежен момент \cite{ref10}. PWM применет на спроводните прекинувачи го поставува ефективниот напон и оттука работната точка на брзината/моментот \cite{ref11}. Без инверторот нема начин да се комутира безчеткеста машина, што е токму причината зошто енергетската електроника --- а не четките --- го дефинира погонот.

\subsection{Како работи инверторското коло}

Енергетскиот степен е трифазен мост: три идентични „гранки“ (полумостови), по една за секоја фаза на моторот, секоја составена од \eng{high-side} прекинувач и \eng{low-side} прекинувач во серија преку еднонасочната собирница, со фазата на моторот поврзана на средната точка. Во овој дизајн секој прекинувач е два IPT015N10N5 MOSFET транзистори поврзани паралелно, па една гранка е четири MOSFET транзистори, а целиот мост е дванаесет.

Во рамките на една гранка, двата прекинувачи работат комплементарно --- никогаш и двата вклучени истовремено, што би ја кратко споило еднонасочната собирница („\eng{shoot-through}“). Вклучувањето на \eng{high-side} прекинувачот ја поврзува фазата со позитивната собирница; вклучувањето на \eng{low-side} прекинувачот ја поврзува со масата; оставањето и на двата исклучени ѝ дозволува на фазата да лебди (се користи за неенергизираната фаза и за \eng{back-EMF} следење) \cite{ref10,ref11}. Чекорењето на трите гранки низ шестчекорната шема ја циркулира струјата низ парови фази по редоследот што произведува ротирачко статорско поле \cite{ref10}. Телесните диоди на MOSFET транзисторите (и индуктивноста на моторот) ја обезбедуваат \eng{freewheeling} патеката за индуктивната фазна струја за време на \eng{dead-time}-то меѓу состојбите на прекинувачите; телесната диода на IPT015N10N5 е карактеризирана токму за оваа намена ($V_{SD}\approx 0.85$\,V, $Q_{rr}\approx 316$\,nC) \cite{ref1}. Спроводната состојба на MOSFET-от со отпорнички карактер ($R_{DS(on)}$ дури 1.5\,m$\Omega$, повторно преполовена со паралелно поврзување) ги држи спроводните загуби ниски при струите од редот на 100\,A \cite{ref1,ref5}.

\subsection{Намена на секоја компонента}

\begin{itemize}
  \item \textbf{Енергетски MOSFET транзистори (IPT015N10N5 $\times 2$ по позиција):} контролираните прекинувачи што ја поврзуваат секоја фаза со собирниците. Два паралелно ја преполовуваат ефективната отпорност во спроводна состојба и ја делат спроводната загуба/топлина; нивниот позитивен температурен коефициент на $R_{DS(on)}$ го прави делењето самостабилизирачко \cite{ref5}.
  \item \textbf{\eng{Gate driver} (UCC27211A-Q1):} ја претвора PWM логиката на контролерот во високострујното 12\,V управување на гејтот потребно за брзо прекинување на MOSFET транзисторите, го поместува нивото на \eng{high-side} сигналот до лебдечкиот јазол на прекинување и спроведува UVLO заштита \cite{ref3}. \eng{Gate driver}-и постојат токму затоа што логичкиот излез (често 3.3\,V, $\le 1$\,A) не може ниту целосно да спроведе енергетски MOSFET, ниту да ја испорача врвната струја за брзо прекинување, и затоа што гејтот на \eng{high-side} уредот мора да биде референциран кон јазол што се ниша до напонот на собирницата \cite{ref3,ref5}.
  \item \textbf{\eng{Bootstrap} кондензатор:} лебдечкиот резервоар на енергија што го напојува \eng{high-side} драјверот додека јазолот за прекинување е висок (в.\,\S\,\ref{sec:boottheory}).
  \item \textbf{\eng{Bootstrap} диода (внатрешна):} еднонасочен вентил што го надополнува \eng{bootstrap} кондензаторот од $V_{DD}$ кога јазолот за прекинување е низок и го блокира напонот на собирницата кога е висок \cite{ref3,ref6}.
  \item \textbf{Отпорници на гејтот:} ја поставуваат/ограничуваат струјата на гејтот, го наштелуваат $dV/dt$ заради EMI и ги пригушуваат/балансираат паралелните гејтови \cite{ref5}.
  \item \textbf{\eng{Pull-down} отпорници гејт--извор:} ги држат MOSFET транзисторите исклучени при вклучување на напојувањето и додаваат $dV/dt$ имунитет на вклучување \cite{ref5}.
  \item \textbf{Кондензатори за раздвојување:} ги напојуваат брзите струјни импулси на гејтот локално, така што 12\,V собирницата и преднапонот на драјверот не се рушат при прекинување \cite{ref3}.
  \item \textbf{Капацитивност на DC-link:} ја вкрутува собирницата против големата бранувачка струја на полумостот и го ограничува надскокот при комутација во рамки на оценката од 100\,V на уредот.
\end{itemize}

\subsection{Bootstrap колото: работа, потреба и улогата на Miller plateau}
\label{sec:boottheory}

\paragraph{Зошто полумост само со N-канални транзистори бара bootstrap напојување.} N-каналните MOSFET транзистори се претпочитани за полумостот бидејќи, за дадена површина на чипот и напон, имаат многу пониска $R_{DS(on)}$ од P-каналните уреди. Но N-канален MOSFET се вклучува само кога неговиот \textbf{гејт е неколку волти над неговиот извор}. Кај \eng{low-side} уредот, изворот е на маса, па 12\,V управување референцирано кон маса работи директно. Кај \eng{high-side} уредот, изворот е јазолот за прекинување, кој се крева до полниот напон на собирницата кога тој уред е вклучен; за да остане вклучен, гејтот мора да се управува $\approx 12$\,V \emph{над собирницата} --- напон што не постои никаде во систем со една собирница \cite{ref3,ref5}. \eng{Bootstrap} колото го создава ова лебдечко \eng{high-side} напојување евтино, користејќи само диода и кондензатор \cite{ref6}.

\paragraph{Како работи.} Кога \eng{low-side} прекинувачот е вклучен, јазолот за прекинување (и пинот HS) се влече кон маса; \eng{bootstrap} кондензаторот, поврзан меѓу HB и HS, се полни од 12\,V $V_{DD}$ собирницата низ \eng{bootstrap} диодата до $\approx V_{DD}-V_F \approx 11$\,V \cite{ref3,ref6}. Кога потоа \eng{high-side}-от се вклучува, јазолот за прекинување се крева до напонот на собирницата; \eng{bootstrap} диодата станува обратно поларизирана (го блокира напонот на собирницата), а наполнетиот кондензатор „се вози нагоре“ на јазолот за прекинување, држејќи го пинот HB $\approx 11$\,V над HS и така обезбедувајќи го лебдечкиот преднапон што го држи \eng{high-side} MOSFET-от целосно спроводен \cite{ref3,ref6}. Кондензаторот повторно се дополнува во секој циклус кога \eng{low-side}-от следно се вклучува --- поради што е потребно минимално време на спроводност на \eng{low-side}-от (освежување), и поради што \S\,\ref{sec:cboot} мораше да ги провери и опаѓањето и полнењето \cite{ref6,ref8}.

\paragraph{Каде влегува Miller plateau.} За време на секое вклучување, напонот на гејтот расте, паузира на рамно „\eng{Miller plateau}“ ($\approx 4.4$\,V овде \cite{ref1}), потоа повторно расте. Паузата настанува бидејќи, штом каналот ја носи полната струја на оптоварување, целата струја за управување на гејтот се пренасочува во празнење на капацитивноста гејт--одвод $C_{GD}$, така што напонот на одводот може да се ниша --- за време на овој интервал $V_{GS}$ е стегнат и се движи само напонот на одводот \cite{ref5}. Интервалот на платото (заедно со интервалот од прагот до платото) е токму прозорецот во кој уредот истовремено издржува висок напон и висока струја, па тоа е местото каде што се генерира загубата при прекинување \cite{ref5}. \eng{Miller} полнежот $Q_{gd}$ што мора да се протурка низ платото е удвоен со паралелното поврзување (овде $\approx 102$\,nC во најлош случај), а \emph{брзината} на напонската транзиција е поставена од тоа колку струја драјверот може да втисне во гејтот на нивото на платото --- а не од неговата врвна струја при полн $V_{DD}$ \cite{ref5}.

\paragraph{Зошто се потребни големи изворни и понорни струи.} Анализата на Balogh го истакнува тоа експлицитно: најважната карактеристика на еден \eng{gate driver} е неговата способност за извор/понор \emph{околу напонот на \eng{Miller plateau}-то} ($\sim 5$\,V), бидејќи таа струја одредува колку брзо се одвива напонската транзиција на одводот и оттука загубата при прекинување \cite{ref5}. Голема \textbf{изворна} струја (3.7\,A) брзо го управува гејтот низ платото при вклучување, рушејќи го $V_{DS}$ за $\approx 28$\,ns (\S\,\ref{sec:gateR}) и минимизирајќи ја загубата при вклучување \cite{ref3,ref5}. Голема \textbf{понорна} струја (4.5\,A) го прави истото при исклучување и, што е клучно, го држи гејтот на \emph{исклучениот} уред цврсто на маса додека спротивниот прекинувач го ниша споделениот јазол: силниот понор ја надвладува \eng{Miller} струјата инјектирана низ $C_{GD}$, спречувајќи $dV/dt$-индуцирано (\eng{Miller}) вклучување и резултирачкиот \eng{shoot-through} \cite{ref3,ref5}. Со два паралелни MOSFET транзистори со висок $Q_g$ по позиција, оваа висока струјна способност е тоа што прави еден канал на драјверот да биде доволен --- и затоа TI ја пласира струјата на уредот специфично за „управување на големи енергетски MOSFET транзистори \ldots{} низ \eng{Miller plateau}-то“ \cite{ref3}.

\subsection{Архитектура на UCC27211A-Q1}

TI го дели уредот на пет функционални блокови: влезни степени, UVLO заштита, поместување на ниво (\eng{level shift}), \eng{boot} диода и излезни драјверски степени \cite{ref3}.
\begin{itemize}
  \item \textbf{Влезни степени.} Независни HI и LI влезови со $\approx 68$\,k$\Omega$ \eng{pull-down}, $\approx 4$\,pF влезна капацитивност и TTL-компатибилни прагови (растечки 2.3\,V, опаѓачки 1.6\,V) со широка хистереза за имунитет на шум; влезовите поднесуваат $-10$\,V до $+20$\,V и се независни од $V_{DD}$, па прифаќаат и логика од микроконтролер и сигнали од аналоген контролер \cite{ref3}.
  \item \textbf{UVLO заштита.} Посебни монитори на $V_{DD}$ и на лебдечкото $V_{HB}$--$V_{HS}$ напојување ги принудуваат излезите ниски додека преднапонот не е валиден. Прагот на $V_{DD}$ е 7\,V (хистереза 0.5\,V), а симетричниот \eng{high-side} праг е 6.7\,V (хистереза 1.1\,V), така што \eng{high-side}-от и \eng{low-side}-от се исклучуваат заедно \cite{ref3}.
  \item \textbf{Поместување на ниво (\eng{level shift}).} Ја преведува командата на \eng{high-side} влезот референцирана кон маса нагоре до лебдечкиот \eng{high-side} излезен степен референциран кон HS, со тесно ($\approx 4$\,ns типично) усогласување на доцнењето меѓу каналите за да се избегне временска нерамнотежа во полумостот \cite{ref3}.
  \item \textbf{\eng{Boot} диода.} 120\,V \eng{bootstrap} диодата е интегрирана (анода $V_{DD}$, катода HB), елиминирајќи надворешна диода и освежувајќи го $C_{BOOT}$ во секој циклус кога HS оди ниско \cite{ref3}.
  \item \textbf{Излезни степени.} \eng{Low-side} степенот е референциран $V_{DD}\rightarrow V_{SS}$, а \eng{high-side} степенот $V_{HB}\rightarrow V_{HS}$; двата обезбедуваат 3.7\,A извор / 4.5\,A понор со брзи (7.2\,ns пораст / 5.5\,ns опаѓање во 1\,nF) излези со ниска импеданса и доцнење на простирање од 20\,ns \cite{ref3}.
\end{itemize}

\paragraph{Проектни равенки дадени во каталошкиот лист на UCC27211A-Q1.} Дисипацијата на драјверот се дели на еднонасочен и прекинувачки дел \cite{ref3}:
\begin{equation}
P_{DISS} = P_{DC} + P_{SW}, \qquad P_{DC} = I_Q \times V_{DD}
\end{equation}
со $I_Q < 0.17$\,mA, па $P_{DC}$ е занемарлив \cite{ref3}. Прекинувачкиот дел се изведува од енергијата за полнење на гејтот како еквивалентен кондензатор \cite{ref3}:
\begin{equation}
E_G = \tfrac{1}{2}\,C_{LOAD}\,V_{DD}^{2}, \qquad P_G = C_{LOAD}\,V_{DD}^{2}\,f_{SW} = Q_G\,V_{DD}\,f_{SW}
\end{equation}
каде што гејтот на MOSFET-от е претставен преку неговиот полнеж низ $Q_G = C_{LOAD}\,V_{DD}$ \cite{ref3}. Половина од $P_G$ се дисипира при вклучување, а половина при исклучување; со надворешни отпорници на гејтот оваа загуба се дели меѓу драјверот и отпорниците \cite{ref3}. Истиот израз се појавува кај Balogh како загуба за управување на гејтот, $P_{GATE}=V_{DRV}\,Q_G\,f_{DRV}$, чиј член $Q_G f_{DRV}$ е просечната струја на преднапон на гејтот \cite{ref5}. За овој дизајн, моќноста за управување на гејтот по драјвер (двата прекинувачи на една гранка, 12\,V, 40\,kHz, $\approx 465$\,nC секој) е приближно:
\begin{equation}
P_G \approx 2 \times Q_{g,\text{поз}}\,V_{DD}\,f_{SW} = 2 \times 465\,\mathrm{nC} \times 12\,\mathrm{V} \times 40\,\mathrm{kHz} \approx 0.45\,\mathrm{W}\ \text{по драјвер,}
\end{equation}
поделена меѓу драјверот и осумте отпорници на гејтот на гранката --- удобно во рамки на дисипацијата на SO-PowerPAD куќиштето кога е изведено според упатствата на TI (драјвер близу MOSFET транзисторите, раздвојување директно кај пиновите, тесни кола на гејтот) \cite{ref3}. Барањето за врвна струја е сопствениот проектен критериум на каталошкиот лист: драјверот мора да го испорача \eng{Miller} полнежот во рамки на целното време на транзиција, $I_{PEAK} = Q_{GD}/dt$, и способноста од 3.7\,A е проверена наспроти ова во \S\,\ref{sec:gateR} \cite{ref3}.

%===============================================================
\section{Заклучок}
\label{sec:concl}

Infineon IPT015N10N5 и Texas Instruments UCC27211A-Q1 формираат компатибилен пар за BLDC инвертор од 3--5\,kW, 40\,kHz, напојуван од 24-V / 48-V батерии чиј напон варира од 20 до 30\,V, односно од 40 до 60\,V. При најлошиот напон на собирницата (60\,V, полна 48\,V батерија), 100\,V MOSFET-от дава DC маргина на блокирање од $\approx 1.67\times$ со околу 40\,V резерва за надскокот при исклучување, додека при најнискиот напон (20\,V, празна 24\,V батерија) најлошата струја достигнува $\approx 150$\,A на собирницата ($\approx 75$\,A по уред) --- сè уште далеку под континуираната оценка од $2\times 250$\,A на паралелниот пар \cite{ref1}. Неговиот низок полнеж на гејтот и отпорничка спроводна состојба се погодни за работа со висока фреквенција \cite{ref1}. 120\,V, 3.7\,A/4.5\,A драјверот, напојуван на 12\,V, е удобно во рамки на својот работен прозорец, ја интегрира \eng{bootstrap} диодата и --- единствената ставка што навистина бараше проверка --- има доволно изворна/понорна струја за да го прекине \emph{паралелниот} полнеж на гејтот за под 30\,ns при умерени $\approx 2$\,V/ns (на 60\,V), далеку под неговата граница од 50\,V/ns \cite{ref3}. Главната последица од паралелното поврзување на два уреди со висок $Q_g$ е дека \eng{bootstrap} и $V_{DD}$ кондензаторите за раздвојување мора да бидат $\approx 1$--2\,$\mu$F наместо номиналните 0.1\,$\mu$F од каталошкиот лист, резултат што директно произлегува од равенките за опаѓање на полнежот \cite{ref3,ref6,ref7}. Целосниот сет помошни компоненти за управување на гејтот е: \textbf{$C_{BOOT}=2.2\,\mu$F, $C_{VDD}=4.7\,\mu$F + 0.1\,$\mu$F, асиметрична мрежа на гејтот од $R_{on}=3.3\,\Omega$ (0805, $1/8$\,W) со паралелен $R_{off}=1.0\,\Omega$ (0805, $1/8$\,W) + \eng{Schottky} OFF диода (Nexperia PMEG4030ER) по уред за меко вклучување и брзо исклучување, и $R_{GS}=10$\,k$\Omega$ (0603, $1/10$\,W) по уред}, со опционална 100\,V надворешна \eng{boot} диода. Бидејќи изворот е батерија со значителна импеданса, нискоиндуктивна \eng{DC-link} банка (оценета $\ge 80$--100\,V) близу полумостот е особено важна. Термичкото управување со MOSFET транзисторите (ладење за искористување на патеката спој--куќиште од 0.4\,K/W; во најлош случај $\approx 8.4$\,W по уред при празна 24\,V батерија) и нискоиндуктивен распоред остануваат задачи на ниво на плоча, надвор од опсегот на вредности на компоненти на оваа студија, но се суштински за остварување на пресметаните прекинувачки перформанси \cite{ref1,ref3}.

%===============================================================
\begin{thebibliography}{99}
\setlength{\itemsep}{3pt}

\bibitem{ref1} Infineon Technologies, \emph{IPT015N10N5 --- OptiMOS\textsuperscript{\texttrademark}\,5 Power-Transistor, 100\,V}, Final Data Sheet Rev.\,2.4, 2023. \url{https://www.infineon.com/assets/row/public/documents/24/49/infineon-ipt015n10n5-datasheet-en.pdf}

\bibitem{ref2} Infineon Technologies, \emph{IPT015N10N5 product page (parametric summary)}. \url{https://www.infineon.com/part/IPT015N10N5}

\bibitem{ref3} Texas Instruments, \emph{UCC27211A-Q1 Automotive 120\,V, 3.7\,A/4.5\,A Half-Bridge Driver with 8\,V UVLO}, Datasheet SLUSCG0C, Dec.\,2015 --- rev.\,Sep.\,2025. \url{https://www.ti.com/lit/ds/symlink/ucc27211a-q1.pdf}

\bibitem{ref4} Texas Instruments, \emph{UCC27211A product folder}. \url{https://www.ti.com/product/UCC27211A}

\bibitem{ref5} L.\,Balogh, \emph{Fundamentals of MOSFET and IGBT Gate Driver Circuits}, Texas Instruments Application Report SLUA618A (rev.\,of SLUP169), 2017--2018. \url{https://www.ti.com/lit/ml/slua618/slua618.pdf}

\bibitem{ref6} Infineon Technologies, \emph{Using Monolithic High-Voltage Gate Drivers} (bootstrap sizing, gate-resistor and half-bridge guidance; successor to IR app notes AN-978/AN-1123), Application Note. \url{https://www.infineon.com/assets/row/public/documents/24/42/infineon-using-monolithic-high-voltage-gate-drivers-applicationnotes-en.pdf}

\bibitem{ref7} Infineon Technologies, \emph{A typical calculation of a bootstrap capacitor's value}, Knowledge Base Article. \url{https://community.infineon.com/t5/Knowledge-Base-Articles/A-typical-calculation-of-a-bootstrap-capacitor-s-value/ta-p/739554}

\bibitem{ref8} Infineon Technologies, \emph{Bootstrap circuits --- overview and selection guide}, Knowledge Base Article. \url{https://community.infineon.com/t5/Knowledge-Base-Articles/Bootstrap-circuits-overview-and-selection-guide/ta-p/720580}

\bibitem{ref9} Infineon Technologies, \emph{Calculating maximum bootstrap capacitance for integrated-diode gate drivers}, Knowledge Base Article. \url{https://community.infineon.com/t5/Knowledge-Base-Articles/Calculating-maximum-bootstrap-capacitance-for-integrated-diode-gate-drivers/ta-p/1015593}

\bibitem{ref10} onsemi, \emph{Trapezoidal Control of BLDC Motors}, technical blog. \url{https://www.onsemi.com/company/news-media/blog/industrial/en-us/trapezoidal-control-of-bldc-motors}

\bibitem{ref11} NXP Semiconductors, \emph{3-phase BLDC Motor Control with Sensorless Back-EMF Zero Crossing Detection}, Application Note AN1913. \url{https://www.nxp.com/docs/en/application-note/AN1913.pdf}

\bibitem{ref12} Nexperia, \emph{PMEG4030ER --- 40\,V, 3\,A low $V_{F}$ Schottky barrier rectifier}, Data Sheet, 2023. \url{https://assets.nexperia.com/documents/data-sheet/PMEG4030ER.pdf}

\bibitem{ref13} Nexperia, \emph{Power MOSFET gate driver fundamentals}, Application Note AN90059, 2025. \url{https://assets.nexperia.com/documents/application-note/AN90059.pdf}

\bibitem{ref14} Toshiba Electronic Devices \& Storage, \emph{Impacts of the dv/dt Rate on MOSFETs}, Application Note, 2018. \url{https://toshiba.semicon-storage.com/info/docget.jsp?did=59464}

\end{thebibliography}

\vspace{1em}
\noindent\rule{\textwidth}{0.4pt}\\[4pt]
{\small \emph{Забелешки за методот:} Сите параметри на уредите и управувачките равенки се преземени од примарните документи на производителите наведени погоре. Нумеричките резултати се пресметани за најлошиот случај (максималните каталошки вредности) таму каде што тие влијаат на безбедносна или димензиона маргина (на пр.\,максимален полнеж на гејтот за димензионирање на кондензаторот). Равенките за \eng{bootstrap} кондензаторот, загубата за управување на гејтот и \eng{Miller plateau}-то се репродуцирани од \cite{ref3}, \cite{ref5} и \cite{ref6,ref7,ref8,ref9}; дизајнот ја заменува работната точка на оваа апликација (12\,V управување, 40\,kHz, два паралелни IPT015N10N5 по прекинувач) во тие објавени релации.}

\end{document}
